多 AMR 仓库调度的运行效果由资源配置与约束紧性共同决定。本部分围绕多 AMR 仓库调度系统的性能退化原因开展灵敏度分析。在本项目中，仓库拓扑、任务规则和动态事件机制给定后，调度系统的主要可控因素集中在 \textbf{AMR 数量、任务负荷和通道通行能力}。AMR 数量决定可用运输资源，任务负荷决定资源需求强度，通道容量决定局部路网约束。三者均可由工程措施直接改变，并且会同时影响任务分配、车辆排序、路径占用和动态重排结果。

在实际运营中，我们需要知道真正制约系统的资源瓶颈。若 \textbf{增加 AMR} 后完工时间改善明显，扩车具有直接价值；若任务量增加后等待和重排快速上升，说明系统已接近 \textbf{任务负荷边界}；若少数通道长期高占用并引发冲突，说明应优先关注 \textbf{局部通行能力}。灵敏度分析的意义就在于把这些现象拆成可比较的实验结果，支撑后续投入选择。

因此，我们以当前仓库布局、任务集合和动态事件为基准，依次进行 \textbf{基线校验}、\textbf{AMR 数量边际分析}、\textbf{任务负荷边界识别}、\textbf{热点通道容量检验} 和 \textbf{工程干预对比}。最终结论服务于三个直接决策：是否继续扩充 AMR，是否控制单位周期内的任务释放量，是否优先改造瓶颈区域。
